\section{Conclusion}
\label{ch:conclusion}

% =============================================================================
% SKELETON. Written last, once Section~\ref{ch:results} exists.
% BUDGET: 650 words of prose. Sections 7.1-7.3 wait on results; 7.4 is drafted
% below and is already at budget.
% =============================================================================

\subsection{Summary of Findings}
\label{sec:summary}

% BUDGET: 250 words.
% CONTENT REQUIRED:
% - One short paragraph per research question: answer it in a sentence, qualify it in
%   another. Do not restate Section~\ref{ch:results}.
% - Be explicit about which findings are predictions (BNPL on, parameters held fixed)
%   and which are calibrated agreement (the baseline).
% - The falsified non-linearity hypothesis belongs here as prominently as the successes,
%   and it belongs in the abstract too.

\subsection{Contributions}
\label{sec:contributions}

% BUDGET: 200 words, in descending order of defensibility.
% 1. An ABM containing a regulated, bureau-informed lender class and an unregulated,
%    bureau-invisible one, over a survey-derived household population. This is the gap
%    of Section~\ref{sec:lit-gap} --- state the novelty claim no more strongly than the
%    systematic search supports.
% 2. A survey-grounded synthetic population of South African households with complete
%    2017-Rand balance sheets and its validation record (Section~\ref{ch:population}),
%    reusable independently of this model.
% 3. A quantification, internal to the model, of the cost of the BNPL bureau-reporting
%    exemption, from the bureau-visibility contrast --- including the finding that it is
%    close to zero, which is the opposite of what the design anticipated.

\subsection{Policy Implications}
\label{sec:policy}

% BUDGET: 200 words.
% CONTENT REQUIRED:
% - Follow the intervention ranking of Section~\ref{sec:results-rq3}: the three
%   instruments regulators are proposing do little here, and the one that works is
%   proposed nowhere. That asymmetry is the policy contribution.
% - Keep the hypothetical concurrent-facility cap clearly separated from the three
%   instruments that exist in other jurisdictions.
% - Honour Section~\ref{sec:lim-reading}: recommend the direction of a policy, never a
%   threshold value.

\subsection{Future Work}
\label{sec:future-work}

Three extensions follow from the limitations of Section~\ref{ch:discussion}, ordered by the ratio
of what they would add to what they would cost.

\textbf{Contagion channels.} The model produces correlated individual defaults rather than systemic
cascades, having no mechanism by which one household's default affects another's income
(Section~\ref{sec:lim-structural}). This is the direct cause of the central negative result, so it
is also the most valuable thing to relax. Three candidates would do it: a traditional lender with a
balance sheet that can be impaired and can ration credit in response; feedback from aggregate
default to employment or income; and Cardaci's expenditure-cascade mechanism
\cite{cardaci2018inequality}, in which peer influence acts on consumption rather than only on
adoption. Any of the three would let a successor model ask the systemic question this one cannot.

\textbf{Bounding the adoption parameters.} The most valuable improvement available without new data
is to constrain $q_{\text{base}}$ using the \ac{CFPB} stacking shares as a calibration target and
not only as an after-the-fact check (Section~\ref{sec:lim-transfer}). This converts a free
parameter into a partly calibrated one at the cost of one additional calibration loop, with no
structural change.

\textbf{Behavioural evidence for the South African market.} The transfer problem cannot be solved
from within a modelling exercise. Eliciting present bias, minimum-payment anchoring or \ac{BNPL}
social-approval effects in a South African sample would replace the model's least defensible
imported parameters with domestic estimates.

